% New in Phase R1 (D28). Carries what Proposition 1 leaves out: the residual assumptions in the
% lifted coordinates, the constants L_g and kappa, the uniform range R, the cost bound, and the
% proof of both lines of the proposition.
\section{Assumptions and Proof for the Selection Bound, Eq.~\eqref{eq:selection_bound}}
\label{app:A6_selection_bound}

The selection bound of Eq.~\eqref{eq:selection_bound} is stated on the memory state, which is the
form the reported instance uses. The proof is written in the lifted
coordinates of Appendix~\ref{app:A1_estimator}, so that it covers the learned
variant as well. Set $\xi_t=E_{\phi}(s_t)$ and write $g_{\psi}=C\circ D_{\psi}$
for the attribute readout. The identity lifting recovers the main-text
statement, with $\xi_t=s_t$ and $g_{\psi}=C$.

For a candidate $c$, let $s_h^c$ and $y_h^c$ be the memory state and the
attribute after $h$ further interactions that hold $c$, starting from
$s_0^c=s_t$, and set $\xi_h^c=E_{\phi}(s_h^c)$. The residuals of the fitted
model along such a rollout are what the bound is stated in terms of,
\begin{equation}
    \begin{aligned}
        \xi_{h+1}^c
        &=K\xi_h^c+Bc+b+\eta_h^c,
        &
        \|\eta_h^c\|_2
        &\leq\varepsilon_{\mathrm{dyn}},\\
        y_h^c
        &=g_{\psi}(\xi_h^c)+\nu_h^c,
        &
        |\nu_h^c|
        &\leq\varepsilon_{\mathrm{rec}}.
    \end{aligned}
    \label{eq:residual_assumptions}
\end{equation}
Here $\eta_h^c$ is the one-step residual of the fitted transition and
$\nu_h^c$ is the residual left in the readout. Both bounds are assumed to hold
over the candidate rollouts, and both absorb imperfect finite-memory closure,
model approximation and generation variability. Assume further that
$\|K\|_2\leq\kappa$, and that $g_{\psi}$ is $L_g$-Lipschitz on the region the
rollouts visit. For stochastic generation the statements hold on the event
that these conditions are met.

The predicted and the actual rollout start from the same encoding, so the gap
between them is the accumulated one-step residual carried through the readout,
\begin{equation}
    \left|
        y_h^c-\widehat{y}_{t+h\mid t}(c)
    \right|
    \leq
    \varepsilon_{\mathrm{rec}}
    +
    L_g\varepsilon_{\mathrm{dyn}}
    \sum_{j=0}^{h-1}\kappa^j
    \;\triangleq\;\Delta_h .
    \label{eq:prediction_bound_full}
\end{equation}
Here $\Delta_h$ is the $h$-step attribute error bound of
the selection bound of Eq.~\eqref{eq:selection_bound}. It trades against the horizon: a longer
horizon covers more distant consequences and asks the dynamics to hold over
more steps. A contraction, $\kappa<1$, would bound the geometric sum by
$1/(1-\kappa)$, and no contraction is assumed for the fitted model.

The cost of a candidate is a weighted sum of squared deviations, so a bound on
each deviation gives a bound on the cost. Let $R$ bound how far any attribute,
predicted or actual, sits from the reference across candidates and rollout
steps,
\begin{equation}
    \left|
        J_t(c)-\widehat{J}_t(c)
    \right|
    \leq
    2R\sum_{h=1}^{H}w_h\Delta_h
    \;\triangleq\;\delta_J .
    \label{eq:cost_bound}
\end{equation}
Here $\delta_J$ bounds the gap between the predicted and the actual
held-command cost of one candidate. The factor $2R$ comes from the identity
$a^2-b^2=(a+b)(a-b)$ applied to each deviation pair.

Let $c_t^{\mathrm{opt}}\in\operatorname*{arg\,min}_{c\in\mathcal{C}_t}J_t(c)$.
Two applications of the cost bound around the definition of $c_t^{\star}$
close the argument,
\begin{equation}
    \begin{aligned}
        J_t(c_t^{\star})
        &\leq
        \widehat{J}_t(c_t^{\star})+\delta_J\\
        &\leq
        \widehat{J}_t(c_t^{\mathrm{opt}})+\delta_J\\
        &\leq
        \min_{c\in\mathcal{C}_t}J_t(c)+2\delta_J.
    \end{aligned}
    \label{eq:selection_bound_proof}
\end{equation}
The middle line holds because $c_t^{\star}$ minimizes the predicted cost. This
is the second line of the selection bound of Eq.~\eqref{eq:selection_bound}. It concerns the
current held-command comparison, not the cumulative cost of the replanned
loop, and it assumes nothing about the candidate set beyond its being finite
and nonempty.
